\documentclass[10pt,a4paper]{article}
\usepackage{rpgcover}
\usepackage{illustrated-panel}
\usepackage{tocloft}
\babelprovide[import]{french}
\selectlanguage{french}
\RPGSetRunningTitle{Au dernier lancer}
\setlength{\cftsecnumwidth}{4.5em}
\setlength{\cftsecindent}{0em}
\renewcommand{\cftsecfont}{\normalsize\bfseries}
\renewcommand{\cftsecpagefont}{\normalsize\bfseries}
\setlength{\cftbeforesecskip}{3pt}
\renewcommand{\cftsecfont}{\bfseries}
\renewcommand{\cftsecpagefont}{\bfseries}
\setlength{\parskip}{2pt plus .5pt minus .5pt}
\linespread{1.025}\selectfont
\hypersetup{pdftitle={Au dernier lancer — édition française},pdfauthor={Jacopo Biscella}}
% French-only chapter pagination; shared theme remains untouched.
\renewcommand{\RPGChapter}[1]{\par\needspace{6\baselineskip}\addvspace{21pt}%
 \refstepcounter{rpgchapter}\phantomsection\addcontentsline{toc}{section}{\protect\numberline{\therpgchapter.}#1}%
 \begin{center}{\fontsize{19}{23}\selectfont\bfseries \therpgchapter\,\textperiodcentered\,\MakeUppercase{#1}\par}%
 \vspace{7pt}{\color{rpgline}\rule{.72\linewidth}{.5pt}}\end{center}\vspace{8pt}}
\newcommand{\RPGFrenchContents}{%
\thispagestyle{rpgfrontmatter}\vspace*{3mm}
\begin{center}{\fontsize{26}{30}\selectfont\bfseries SOMMAIRE\par}\vspace{9pt}
{\color{rpgline}\rule{.76\linewidth}{.55pt}}\end{center}\vspace{3mm}
\begingroup\csname @starttoc\endcsname{toc}\endgroup
\vspace{3mm}
\begin{center}\includegraphics[width=.51\textwidth,height=43mm,keepaspectratio]{assets/last-throw-crosshatch-redrawn.png}\end{center}
\vfill
\begin{center}{\small\itshape Un jeu de Jacopo Biscella\par}\vspace{5pt}
{\small Cette œuvre est mise à disposition selon les termes de la licence Creative Commons Attribution — Partage dans les mêmes conditions 4.0 International (CC BY-SA 4.0). Vous pouvez la partager et l'adapter, y compris à des fins commerciales, à condition de créditer l'auteur et de diffuser les œuvres dérivées sous la même licence.\par}\end{center}\clearpage}
\begin{document}
\RPGStartFrontMatter
\RPGSetCoverWordSpace{.78em}
\RPGFrontCover{AU DERNIER LANCER}{CHOIX, DESTIN, CONSÉQUENCES}{Jacopo Biscella}{assets/last-throw-crosshatch-redrawn.png}
\RPGStartContents
\RPGFrenchContents
\RPGStartMainMatter
\RPGChapter{Le Dé}

Chaque fois qu'un Protagoniste entreprend une action dont l'issue est incertaine, il lance des dés à six faces (d6). Le nombre de dés lancés et la manière d'interpréter le résultat dépendent de son niveau dans la Caractéristique utilisée.

\section*{Niveaux et jets}
\begin{rpgtable}{Y|Y|Y}
\rpghead{Niveau} & \rpghead{Jet} & \rpghead{Signification} \\ \hline
Novice & 2d6, garder le plus faible & Hors de son domaine \\ \hline
Compétent & 1d6 & Aptitude fonctionnelle \\ \hline
Expert & 2d6, garder le plus élevé & Formé, maître de sa discipline \\ \hline
\end{rpgtable}

Une \textbf{Compétence} pertinente améliore le jet d'un degré. Une \textbf{pénalité de situation} (blessure, obscurité, fatigue) le dégrade d'un degré. Les deux s'annulent mutuellement.

\begin{rpgtable}{Y|Y}
\rpghead{Situation} & \rpghead{Jet effectif} \\ \hline
Novice + pénalité & Faible : 3d6, garder le plus faible \\ \hline
Novice & 2d6, garder le plus faible \\ \hline
Novice + compétence / Compétent & 1d6 \\ \hline
Compétent + compétence / Expert & 2d6, garder le plus élevé \\ \hline
Expert + compétence & Exceptionnel : 3d6, garder le plus élevé \\ \hline
\end{rpgtable}

\section*{Seuil de difficulté}
Le résultat du dé est comparé à un \textbf{Seuil de difficulté (SD)} : la valeur minimale à atteindre pour réussir. Le Narrateur fixe ce seuil selon la difficulté de l'action.
\begin{rpgtable}{Y|Y}
\rpghead{SD} & \rpghead{Difficulté} \\ \hline
3 & Simple \\ \hline
4 & Ordinaire \\ \hline
5 & Difficile \\ \hline
6 & Extrême \\ \hline
Critique & Miraculeuse — exige deux 6 \\ \hline
\end{rpgtable}
Le Narrateur interprète librement la qualité de la réussite ou de l'échec : un 6 face à un SD de 4 est brillant ; un 4 face à ce même seuil est une réussite de justesse.

\needspace{13\baselineskip}
\section*{Le Critique}
Si le jet donne \textbf{deux 6 ou davantage}, c'est un \textbf{Critique} : un résultat extraordinaire. Un Critique exige au moins deux 6 sur des dés valables pour le Critique. Les dés ajoutés par un désavantage ne comptent pas ; les dés bonus comptent. Sans bonus, seuls les niveaux Expert et Exceptionnel peuvent produire un Critique.

Certaines actions — enfoncer une porte de fer à mains nues, briser un antique enchantement, accomplir un miracle — \textbf{exigent explicitement un Critique}. Seuls les personnages disposant d'une Compétence pertinente peuvent les tenter.

\textbf{Règle d'Épuisement} : après l'échec d'une action exigeant un Critique, une nouvelle tentative n'est possible que si les circonstances changent de manière significative : repos adéquat, nouvelle approche, aide extérieure ou changement de situation. Le Narrateur décide si le changement est suffisant.

\section*{Probabilités en un coup d'œil}
\begin{rpgtable}{Y|Y|Y|Y}
\rpghead{Niveau} & \rpghead{Prob. $\geq 4$} & \rpghead{Prob. 6} & \rpghead{Critique} \\ \hline
Faible (3d6$\downarrow$) & $\sim$12\% & $\sim$0,5\% & — \\ \hline
Novice (2d6$\downarrow$) & $\sim$25\% & $\sim$3\% & — \\ \hline
Compétent (1d6) & $\sim$50\% & $\sim$17\% & — \\ \hline
Expert (2d6$\uparrow$) & $\sim$75\% & $\sim$31\% & $\sim$2,8\% \\ \hline
Exceptionnel (3d6$\uparrow$) & $\sim$87\% & $\sim$42\% & $\sim$7\% \\ \hline
\end{rpgtable}

\RPGChapter{Créer un Protagoniste}
Chaque Protagoniste est défini par une Espèce, une Classe, quatre Caractéristiques, des Compétences, un Trait, un Fardeau et des Points de Vie.

Le joueur dispose de \textbf{10 points} à répartir entre Points de Vie, Caractéristiques et Compétences. La Classe détermine les coûts, les limites et les possibilités.

\section*{Coût des Caractéristiques}
\begin{rpgtable}{Y|Y|Y}
\rpghead{Niveau} & \rpghead{Caract. principale} & \rpghead{Autre Caractéristique} \\ \hline
Novice → Compétent & 1 point & 2 points \\ \hline
Compétent → Expert & 2 points & 4 points \\ \hline
Expert → Exceptionnel & 4 points & 8 points \\ \hline
\end{rpgtable}
Chaque Point de Vie coûte \textbf{1 point}. Chaque Compétence ou Grimoire coûte \textbf{1 point}.

Le \textbf{bonus d'espèce} accorde une réduction d'un point lors de l'acquisition du niveau Compétent dans la Caractéristique de sa famille. Cette étape est gratuite si la Caractéristique est aussi celle de la Classe. La réduction ne s'applique qu'à cette étape.

\section*{Trait}
Une formule qui exprime la nature profonde du personnage. Aucun effet mécanique : elle nourrit les histoires du Narrateur et oriente le joueur.

\textit{« Je n'abandonne personne. » · « Le pouvoir se prend, il ne se demande pas. » · « Je préfère une vérité inconfortable à un mensonge bienveillant. »}

\section*{Fardeau}
Un poids que porte le personnage : regret, dette, secret ou malédiction. Le Narrateur s'en sert pour créer des dilemmes et des complications personnelles.

\textit{« J'ai tué un innocent pendant la guerre. » · « Je dois la vie à quelqu'un que je déteste. » · « Quelqu'un porte mon visage. »}

\RPGChapter{Caractéristiques et Compétences}
Tous les personnages ont quatre Caractéristiques, qui commencent au niveau Novice. Le joueur les améliore en dépensant des points lors de la création.
\begin{rpgtable}{Y|Y}
\rpghead{Caractéristique} & \rpghead{Domaine} \\ \hline
Force & Puissance physique, endurance, combat rapproché \\ \hline
Dextérité & Agilité, réflexes, discrétion, visée, esquive \\ \hline
Savoir & Raisonnement, culture, connaissance, influence sociale \\ \hline
Merveille & Rapport au surnaturel, miracles \\ \hline
\end{rpgtable}
Une Compétence pertinente améliore le jet d'un degré. Le personnage peut acheter une Compétence s'il possède au moins le niveau Compétent dans la Caractéristique associée. Le coût de celle-ci constitue une limite naturelle.

\section*{Compétences de Force}
\textbf{Armes lourdes} — armes à deux mains, haches, marteaux de guerre.

\textbf{Boucliers} — défense active avec un bouclier, aussi utilisé comme arme.

\textbf{Lutte} — combat à mains nues, saisies, immobilisations.

\textbf{Démolition} — destruction d'objets, de portes, de murs, de barricades.

\textbf{Endurance} — résistance à la douleur, la fatigue, les poisons et les conditions extrêmes.

\section*{Compétences de Dextérité}
\textbf{Discrétion} — se déplacer sans être repéré.

\textbf{Acrobatie} — équilibre, sauts et mouvements complexes.

\textbf{Crochetage} — ouvrir serrures et mécanismes.

\textbf{Pièges} — détecter et neutraliser pièges et mécanismes dangereux.

\textbf{Visée} — atteindre précisément des cibles éloignées.

\textbf{Vol à la tire} — dérober des objets discrètement.

\textbf{Équitation} — maîtriser une monture en déplacement et au combat.

\textbf{Survie} — s'orienter à l'instinct, chasser, suivre des traces et camper.

\needspace{10\baselineskip}
\section*{Compétences de Savoir}
\textbf{Diplomatie} — négocier et rechercher des compromis.

\textbf{Rhétorique} — persuader, inspirer, haranguer, parler en public.

\textbf{Tromperie} — manipulations et mensonges.

\textbf{Intimidation} — pression sociale et psychologique.

\textbf{Commerce} — négociations et estimations.

\textbf{Médecine} — diagnostic, soins, stabilisation.

\textbf{Langues} — langues, textes anciens, communication.

\textbf{Histoire} — événements passés, symboles, contextes.

\textbf{Droit} — lois, contrats, administration.

\textbf{Étiquette} — usages de la cour, convenances.

\textbf{Artisanat} — création, réparation et estimation d'objets.

\textbf{Orientation} — cartes, étoiles, itinéraires.

\textbf{Herboristerie} — remèdes naturels, potions végétales, onguents.

\textbf{Tactique} — stratégie militaire, coordination, planification.

\section*{Compétences de Merveille}
\textbf{Canalisation} — concentrer et libérer l'énergie élémentaire.

\textbf{Rituels} — magie cérémonielle, évocations, consécrations, pactes.

\textbf{Alchimie} — transmutation, élixirs, feu alchimique.

\textbf{Occultisme} — nécromancie, démons, malédictions, savoir interdit.

\textbf{Perception arcanique} — percevoir la magie, divination, voir au-delà du voile.

Ces compétences servent également à lancer des Grimoires : chacun nécessite une Compétence de Merveille précise. Le Mystique a accès aux \textbf{Grimoires}, détaillés dans leur propre chapitre.

\needspace{32\baselineskip}
\RPGChapter{Classes}
Chaque Protagoniste appartient à l'une des quatre Classes. Sa Classe fixe la Caractéristique principale (moins chère à améliorer), le nombre minimum et maximum de Points de Vie et les Compétences disponibles.
\begin{rpgtable}{Y|Y|Y|Y|Y}
\rpghead{Classe} & \rpghead{Principale} & \rpghead{PV min.} & \rpghead{PV max.} & \rpghead{Compétences} \\ \hline
Guerrier & Force & 4 & 7 & 5 (Force) \\ \hline
Aventurier & Dextérité & 3 & 5 & 8 (Dextérité) \\ \hline
Érudit & Savoir & 2 & 5 & 14 (Savoir) \\ \hline
Mystique & Merveille & 2 & 4 & 5 (Merveille) + Grimoires \\ \hline
\end{rpgtable}

\section*{Guerrier}
Le rempart du groupe. Peu de subtilités, beaucoup de Points de Vie, accès exclusif aux armures lourdes et aux armes dévastatrices. Sa liste de Compétences est courte, mais sa force est de rester debout et de frapper fort.

\section*{Aventurier}
Polyvalent et adaptable. Des PV moyens et un large éventail de Compétences. Il résout les problèmes avec ingéniosité physique : crochetage, discrétion, attaques précises, esquive des pièges.

\section*{Érudit}
Le savoir est un pouvoir. Sa longue liste de Compétences permet de devenir médecin, diplomate, stratège, marchand ou escroc. Au combat, il soutient le groupe avec Médecine, Tactique, Intimidation et Rhétorique.

\section*{Mystique}
Fragile, mais capable de l'impossible. Ses faibles Points de Vie sont compensés par l'accès aux Grimoires : des sorts que les autres Classes ne peuvent acquérir qu'à grand prix. Chaque Grimoire est un risque, car un échec magique peut toucher tout le groupe.

\section*{Exemples de création}
\begin{rpgnote}
Guerrier nain — Force Compétent : 0 (bonus d'espèce) → Force Expert : +2. 6 PV : 6. Deux Compétences : 2. Total : 10.

Aventurier elfe — Dextérité Compétent : 0 (bonus d'espèce) → Dextérité Expert : +2. 4 PV : 4. Quatre Compétences : 4. Total : 10.

Érudit humain — Savoir Compétent : 0 (bonus d'espèce) → Savoir Expert : +2. 3 PV : 3. Cinq Compétences : 5. Total : 10.

Mystique féerique — Merveille Compétent : 0 (bonus d'espèce) → Merveille Expert : +2. 2 PV : 2. Trois Compétences : 3. Trois Grimoires : 3. Total : 10.
\end{rpgnote}

\RPGEngravedPanel{Exemple de création — Kael, aventurier elfe}{assets/panel-08-engraved.png}{Luca choisit un Aventurier elfe, de la famille de Dextérité. Grâce au bonus d'espèce, sa Dextérité commence gratuitement au niveau Compétent. Il dispose de 10 points.

Dextérité Expert : 2 points (principale). 4 PV : 4 points. Compétences — Discrétion, Crochetage, Pièges et Visée : 4 points. Total : 10.

Trait : « Aucune porte ne me résiste. » Fardeau : « J'ai volé quelque chose que je n'aurais jamais dû toucher. »}

\RPGChapter{Espèces}
Les espèces se répartissent en quatre familles, chacune liée à une Caractéristique. L'espèce accorde une \textbf{réduction d'un point} pour obtenir le niveau Compétent dans la Caractéristique correspondante. Les espèces proposées sont des exemples : Narrateur et joueurs choisissent celles de leur univers ou en inventent d'autres.
\begin{rpgtable}{Y|Y|Y}
\rpghead{Famille} & \rpghead{Caractéristique} & \rpghead{Exemples} \\ \hline
Espèces de Force & Force & Nains, Orcs, Minotaures, etc. \\ \hline
Espèces de Dextérité & Dextérité & Elfes, Halfelins, Gobelins, etc. \\ \hline
Espèces de Savoir & Savoir & Humains, Métamorphes, Sirènes, etc. \\ \hline
Espèces de Merveille & Merveille & Fées, Lutins, Démons, Anges, etc. \\ \hline
\end{rpgtable}

La seule règle mécanique de l'espèce est sa famille, qui détermine la réduction. Un Mystique nain bénéficie d'une Force moins coûteuse, mais paie le prix normal pour la Merveille.

\RPGEngravedPanel{Un Mystique nain ?}{assets/panel-09-engraved.png}{Grunthar est un Nain qui a choisi une vocation inattendue. Trapu, barbu, les mains d'un forgeron, il ressent pourtant la magie du fer. Son espèce lui donne gratuitement Force Compétent, mais son cœur appartient à la Merveille. Cela lui coûte cher : 2 points pour Compétent, 4 pour Expert. Cependant, quand il canalise le feu dans son marteau, plus personne ne rit.}

\RPGChapter{Combat}
Pas de tours rigides, d'initiative ou de grille. Le combat est une conversation structurée.

\section*{Déroulement}
\textbf{1.} Le Narrateur décrit la situation et demande : \textit{« Que fais-tu ? »}

\textbf{2.} Chaque Protagoniste annonce son action.

\textbf{3.} Il lance la Caractéristique pertinente contre le SD fixé par le Narrateur.

\textbf{4.} Le Narrateur décrit les conséquences et le monde réagit.

\textbf{5.} On recommence jusqu'à la résolution du conflit : victoire, fuite, reddition ou mort.

Il n'y a ni initiative ni ordre fixe. Le Narrateur donne la parole selon la logique de la scène : proximité, danger immédiat ou urgence des intentions. En situation chaotique, tout le monde peut déclarer son action en même temps, et le conflit se résout dans cette confusion.

\section*{Fuite}
Fuir est une action comme une autre. Lancez la Dextérité contre un SD déterminé par la situation : semer un gobelin lent exige SD 3 ; échapper à un chevalier monté peut exiger SD 5 ou davantage.

\section*{Jets opposés}
Lorsque deux personnages s'opposent directement, chacun lance les dés. Le résultat le plus élevé l'emporte. En cas d'égalité, la situation reste bloquée sur le plan narratif.

\section*{Aider un autre Protagoniste}
Un Protagoniste peut aider un compagnon. L'effet dépend de sa propre compétence pour l'action concernée.

\textbf{Aide compétente} (l'allié lancerait au niveau Expert ou supérieur pour le même test) : ajoute \textbf{+1d6} à la réserve du compagnon.

\textbf{Aide incompétente} (l'allié lancerait à un niveau inférieur à Expert pour le même test) : ajoute \textbf{+1d6}, mais le \textbf{SD augmente de 1}. Si le SD était déjà de 6, il devient Critique (deux 6). Si le jet exigeait déjà deux 6, il en exige trois, et ainsi de suite.

\textit{L'incompétent qui « aide » avec un ancien grimoire pourrait rendre l'action impossible. Le joueur doit décider si le risque en vaut la peine.}

\section*{Distances}
Le jeu n'utilise ni grille ni mesures exactes. Les distances sont exprimées par quatre notions narratives communes au Narrateur et aux joueurs.
\begin{rpgtable}{Y|Y|Y}
\rpghead{Distance} & \rpghead{Signification} & \rpghead{Exemples} \\ \hline
Au contact & Contact physique, corps à corps & Lutte, dague, saisie \\ \hline
Proche & Quelques pas, accès immédiat & Épée, charge, discussion \\ \hline
À portée & Distance de tir, trop loin pour la mêlée & Arc, Grimoires à distance \\ \hline
Loin & Hors de portée, approche prolongée & Au-delà du tir, fuite \\ \hline
\end{rpgtable}

Se rapprocher d'une catégorie est une action : de Loin à À portée, de À portée à Proche, de Proche à Au contact. Le Narrateur peut exiger un jet pour avancer sous la pression.

\section*{Exemple de combat}
\textit{Kara (Guerrière), Syl (Aventurière) et Venn (Mystique) sont surpris par deux Pillards orcs (Experts, SD 5 pour les toucher).}

\textbf{Narrateur} : \textit{« Deux orcs surgissent des arbres. Ils sont Proches, armés de marteaux. Que faites-vous ? »}

\textbf{Kara} : \textit{« Je charge le premier avec mon marteau. »} Elle lance Force Expert (2d6, garde le plus élevé), obtient 3 et 5 et conserve 5. SD 5 : réussite ! Marteau avec Force = 3 Blessures. L'orc passe de 4 à 1.

\textbf{Syl} : \textit{« Je recule jusqu'à À portée et tire à l'arc. »} Dextérité Expert + Visée (Exceptionnel, 3d6, garder le plus élevé). Résultats : 2, 4, 6. Le 6 suffit ! L'arc inflige 2 Blessures. Le premier orc tombe.

\textbf{Venn} : \textit{« Mur de flammes pour arrêter le second ! »} Le sort coûte 2 Blessures de Fatigue (Venn passe de 3 à 1 PV, état Grave). Il lance Merveille Expert + Canalisation (Exceptionnel, 3d6), mais subit la pénalité Grave ($-2$ degrés, soit 1d6). Il obtient 4 contre SD 5 : échec ! Retour de flamme : les flammes explosent, tout le monde subit 2 Blessures. Venn tombe à $-1$ PV et quitte la scène.

\RPGChapter{Armes et Protection}
Les dégâts dépendent de l'association entre l'arme et la Caractéristique employée. Chaque arme possède une Caractéristique naturelle ; l'employer avec l'autre réduit les dégâts.

\section*{Table des armes}
\begin{rpgtable}{Y|Y|Y|Y}
\rpghead{Arme} & \rpghead{Avec Force} & \rpghead{Avec Dextérité} & \rpghead{Type} \\ \hline
Dague & 1B & 2B & Blessure \\ \hline
Épée & 2B & 2B & Blessure \\ \hline
Marteau de guerre & 3B & 1B & Blessure \\ \hline
Arc & — & 2B & Blessure \\ \hline
Lance & 2B & 2B & Blessure \\ \hline
Hache de bataille & 2B & 1B & Blessure \\ \hline
Masse & 2F & 1F & Fatigue \\ \hline
Bâton & 1F & 1F & Fatigue \\ \hline
Poings (Lutte) & 1F & 1F & Fatigue \\ \hline
\end{rpgtable}

\textbf{B} = Blessure (dégât létal). \textbf{F} = Fatigue (dégât non létal, étourdissement). Un personnage réduit à zéro PV uniquement par la Fatigue s'évanouit sans danger de mort immédiat.

Un \textbf{Critique} au combat inflige \textbf{+1 Blessure} en plus des dégâts de l'arme.

\section*{Protection}
L'armure détermine le \textbf{SD que l'attaquant doit atteindre pour toucher}. Chaque armure possède deux valeurs, contre Force et contre Dextérité.
\begin{rpgtable}{Y|Y|Y|Y}
\rpghead{Armure} & \rpghead{SD vs Force} & \rpghead{SD vs Dext.} & \rpghead{Accès} \\ \hline
Aucune & 3 & 3 & Tous \\ \hline
Légère & 3 & 4 & Toutes les classes \\ \hline
Lourde & 4 & 5 & Guerrier seulement \\ \hline
\end{rpgtable}

\textbf{Restriction} : seul le Guerrier est entraîné à porter une armure lourde. Aventuriers, Érudits et Mystiques ne peuvent pas en équiper, quelles que soient leurs Caractéristiques.

\textbf{Bouclier} : lorsqu'un ennemi vous attaque, le SD de son attaque augmente de 1. Nécessite la Compétence Boucliers et une main libre.

\needspace{12\baselineskip}
\section*{Objets remarquables}
Certains objets possèdent un nom et un effet précis : améliorer d'un degré un jet dans une situation particulière, ou rendre possible une action normalement impossible.

\begin{rpgnote}
Lame de sel — +1 degré contre les créatures surnaturelles.

Clé de l'Ailleurs — ouvre toute serrure, mais chaque utilisation attire une attention indésirable.
\end{rpgnote}

Ce ne sont que des exemples. Le Narrateur crée ces objets selon le même principe : nom évocateur, effet précis et, éventuellement, coût ou conséquence.

\RPGEngravedPanel{Autres objets remarquables}{assets/panel-12-engraved.png}{Couronne de commandement — +1 degré en Rhétorique. Son porteur ne peut pas mentir.

Lanterne des morts — révèle fantômes et esprits invisibles, mais attire les esprits malveillants.

Gants de forge — +1 degré en Démolition. Ils surchauffent après chaque usage : 1F.}

\RPGChapter{Blessures et Récupération}
Les Points de Vie forment une réserve unique qui absorbe les dégâts du combat et les Blessures de Fatigue : le coût de tout effort extraordinaire, qu'il s'agisse de lancer un Grimoire, soulever un rocher ou courir des heures.

\section*{Deux types de blessures, une seule réserve}
\textbf{Blessures de combat} : à zéro PV, le personnage sort de la scène : inconscient, capturé ou potentiellement mourant. La mort survient lorsque les seules Blessures de combat atteignent les PV maximum, selon les règles ci-dessous.

\textbf{Blessures de Fatigue} : à zéro PV uniquement par Fatigue, il s'évanouit sans être en danger de mort. Il se réveille après un court repos. Fatigue et blessures de combat s'additionnent : un Mystique qui dépense 2 PV en sorts puis reçoit 1 Blessure de combat court un véritable danger.

\section*{État du personnage}
\begin{rpgtable}{Y|Y|Y}
\rpghead{PV restants} & \rpghead{État} & \rpghead{Effet} \\ \hline
3+ & En forme & Aucun \\ \hline
$\leq 2$ & Amoché & $-1$ degré à tous les jets \\ \hline
$\leq 1$ & Grave & $-2$ degrés \\ \hline
0 & Hors de scène & Inconscient ou mourant \\ \hline
\end{rpgtable}

\section*{Mort}
Un personnage ne meurt que si les \textbf{Blessures de combat (B) seules} atteignent ou dépassent ses PV maximum. Si ne serait-ce qu'un point de dégât relève de la Fatigue (F), il reste inconscient, mais en vie.

Un personnage à terre avec un mélange de B et F peut être achevé : de nouvelles Blessures de combat portant les B totales à son maximum de PV le tuent. Le groupe dispose d'un délai pour intervenir, mais le coup fatal existe.

\section*{États et conditions}
Poison, feu, cécité, peur et autres conditions n'ont pas de règles spécifiques. Le Narrateur les traite comme des pénalités de situation ($-1$ degré) jusqu'à leur résolution. Une seule mécanique couvre tous ces cas.

\section*{Récupération}
\textbf{Repos court} (quelques heures en sécurité) : récupérer 1 Blessure.

\textbf{Repos long} (une nuit complète en sécurité) : récupérer toutes les Blessures et réinitialiser l'Épuisement des actions Critiques ratées.

\textbf{Soins pendant la scène} : Compétence Médecine + réussite = récupération de 1 Blessure (2 sur un Critique).

\RPGEngravedPanel{La spirale de la chute}{assets/panel-13-engraved.png}{Kara dispose de 6 PV. Un orc lui inflige 2B. Puis elle soulève une porte : 1F. Il lui reste 3 PV : elle est en forme, sans pénalité. L'orc la frappe encore : 2B. Elle tombe à 1 PV, état Grave ($-2$ degrés). Elle doit choisir : fuir, ou déclarer un Effort pour tenter un dernier coup désespéré. Au dernier lancer.}

\RPGChapter{Magie et Grimoires}
Les Grimoires sont des sorts spécifiques que le Mystique acquiert lors de la création. Chacun possède des propriétés définies, comme une carte. La magie n'est pas gratuite : lancer un Grimoire peut coûter des Blessures de Fatigue, et tout échec provoque un Contrecoup pouvant toucher le Mystique, ses alliés ou son environnement.

\section*{Structure d'un Grimoire}
Chaque Grimoire possède six propriétés :

\textbf{Nom} — le nom du sort.

\textbf{Compétence requise} — Canalisation, Rituels, Alchimie, Occultisme ou Perception arcanique.

\textbf{SD} — difficulté de lancement (de 3 à 6, ou Critique).

\textbf{Coût} — Blessures de Fatigue payées au moment du lancement. Les sorts légers peuvent coûter 0 : seul le Contrecoup reste un risque.

\textbf{Effet} — ce qui se produit en cas de réussite.

\textbf{Contrecoup} — ce qui se produit en cas d'échec.

\section*{Exemples de Grimoires}
\textbf{Mur de flammes} — Canalisation, SD 5, coût 2. \textit{Effet} : un mur de feu bloque un passage ou encercle les ennemis. \textit{Contrecoup} : les flammes deviennent incontrôlables ; toutes les personnes présentes subissent 2 Blessures.

\textbf{Toucher guérisseur} — Canalisation, SD 3, coût 1. \textit{Effet} : soigne 2 Blessures d'un allié. \textit{Contrecoup} : les Blessures passent au Mystique au lieu de disparaître.

\textbf{Cercle de protection} — Rituels, SD 4, coût 1. \textit{Effet} : une zone délimitée devient infranchissable pour les ennemis pendant une scène. \textit{Contrecoup} : les alliés y sont également piégés ; nul ne peut sortir.

\textbf{Feu alchimique} — Alchimie, SD 4, coût 1. \textit{Effet} : une fiole de feu explose à l'impact et inflige 2 Blessures dans une zone. \textit{Contrecoup} : la fiole explose dans vos mains ; vous et les personnes Proches subissez 1 Blessure.

\textbf{Voix des morts} — Occultisme, SD 5, coût 2. \textit{Effet} : interroger un défunt qui répond sincèrement à trois questions. \textit{Contrecoup} : l'esprit est hostile et possède temporairement un membre du groupe.

\textbf{Œil du néant} — Perception arcanique, SD 4, coût 1. \textit{Effet} : voir à travers murs et illusions pendant une scène. \textit{Contrecoup} : quelque chose de l'autre côté vous voit aussi et vous a trouvé.

\section*{Créer de nouveaux Grimoires}
Narrateur et joueurs peuvent imaginer d'autres Grimoires en suivant cette structure à six champs. La seule règle : le Contrecoup doit être précis, concret et proportionné à l'Effet. Plus le sort est puissant, plus l'échec est dangereux, et pas uniquement pour le Mystique.

\RPGEngravedPanel{Équilibrer un Grimoire}{assets/panel-15-engraved.png}{SD 3, coût 1 : effet modeste, contrecoup personnel. Seul le Mystique est en danger ; idéal pour un sort utilitaire.

SD 5, coût 2 : effet puissant, contrecoup collectif. Tout le groupe prend un risque et doit décider si cela en vaut la peine.

Critique, coût 3 : miracle. En cas d'échec, catastrophe. À réserver aux situations sans autre issue.}

\section*{Magie pour les autres Classes}
Tout personnage peut acheter des Grimoires s'il possède au moins Merveille Compétent et la Compétence requise. Cependant, améliorer Merveille coûte deux fois plus cher aux non-Mystiques : la magie leur est accessible, mais à un prix très élevé.

\RPGChapter{Progression}
Pas de points d'expérience. À la fin de chaque séance, le groupe répond à trois questions :

\textbf{1.} Avez-vous découvert quelque chose d'important ?

\textbf{2.} Votre Trait ou votre Fardeau vous a-t-il causé des problèmes ?

\textbf{3.} Avez-vous accompli quelque chose qui a changé le cours de l'histoire ?

Chaque réponse positive donne \textbf{1 point de progression}. Ces points se dépensent comme à la création : pour améliorer les Caractéristiques, acquérir des Compétences ou obtenir des Grimoires.

\section*{Déclarer et jouer sa progression}
Lorsqu'il reçoit un point, le joueur doit \textbf{déclarer immédiatement où il l'investit}. Si l'amélioration coûte plusieurs points, ceux-ci sont engagés un par un. Pendant ce temps, le joueur doit \textbf{mettre en scène sa progression} : s'entraîner, chercher un maître, expérimenter ou étudier.

Le Narrateur peut refuser une progression jamais montrée pendant les séances. Entre des aventures séparées par des mois de temps fictif, le temps écoulé peut suffire à la justifier.

\needspace{8\baselineskip}
Au-delà des améliorations mécaniques, un joueur peut :

— Réécrire son \textbf{Trait}, car le personnage a changé.

— Résoudre son \textbf{Fardeau} et en choisir un nouveau.

\RPGEngravedPanel{Grandir au fil du jeu}{assets/panel-16-engraved.png}{Séance 3 : Syl déclare vouloir atteindre le niveau Exceptionnel en Dextérité. Cela coûte 4 points. Elle engage le premier.

Séance 5 : elle cherche un maître archer dans la forêt. Deuxième point.

Séance 8 : elle s'entraîne tous les soirs au camp. Troisième point.

Séance 11 : elle affronte une épreuve impossible et la surmonte. Quatrième point. Dextérité Exceptionnel : elle l'a vécu.}

\RPGChapter{Ennemis}
Un ennemi ne possède pas de fiche complète, seulement quatre éléments :

\textbf{SD} — deux valeurs (Force / Dextérité) indiquant la difficulté pour le toucher. Son niveau de jet en attaque est indiqué séparément.

\textbf{Blessures} — combien de coups il peut encaisser (1 à 3 pour les sbires, 4 à 6 pour les vétérans, 7 ou plus pour les chefs).

\textbf{Arme} — détermine les dégâts selon le même tableau que les joueurs.

\textbf{Comportement} — une phrase guidant ses actions.

\section*{Exemples}
\begin{rpgnote}
Gobelin — SD : F3/D3, Compétent, 2 Blessures, dague (2B Dextérité). « Attaque en bande et fuit s'il est seul. »

Pillard orc — SD : F3/D4, Expert, 4 Blessures, marteau (3B Force), armure légère. « Charge le plus faible. »

Chevalier corrompu — SD : F4/D5, Expert, 6 Blessures, épée (2B), armure lourde. « Recherche le duel. Ignore les faibles. »
\end{rpgnote}

\section*{Groupes}
Des ennemis identiques agissent en un seul bloc. Ils lancent un jet commun, mais infligent +1 Blessure par membre supplémentaire (maximum +3). Chaque Blessure infligée au groupe élimine un membre.

\section*{Capacités spéciales}
Le Comportement d'un ennemi peut mentionner une capacité spéciale en une phrase, sans règles détaillées. Le Narrateur l'applique comme une condition narrative.
\begin{rpgnote}
Troll des marais — SD : F4/D3, Expert, 6 Blessures, griffes (2B Force). « Régénère 1 Blessure par tour ; le feu arrête la régénération. »

Vouivre — SD : F4/D4, Expert, 5 Blessures, morsure (3B Force). « Vole. Son venin inflige 1 Blessure de Fatigue supplémentaire à chaque attaque. »
\end{rpgnote}

\RPGChapter{L'Effort}
\textit{« Je saigne pour y parvenir. »}

Avant un jet, le joueur peut déclarer un \textbf{Effort} : il ajoute \textbf{+1d6} au jet. \textbf{Après} le jet, il subit \textbf{1 Blessure de Fatigue}. Ce dé supplémentaire compte à la fois pour le résultat et pour le Critique.

L'Effort se limite à \textbf{+1d6 par action} et ne se cumule pas. Comme la Blessure intervient après le jet, les pénalités d'état (Amoché, Grave) ne modifient pas le jet en cours, mais seulement les suivants.

\begin{rpgnote}
Un Mystique avec 3 PV a dépensé 1 PV pour un Grimoire. À 2 PV, il est Amoché (-1 niveau). Avec Merveille Expert, il lancerait 1d6. Pour un sort SD 5, il déclare un Effort avant le jet : +1d6, soit 2d6 en gardant le plus élevé. Après le jet, il paie 1F et tombe à 1 PV (Grave). S'il échoue et reçoit un coup, il sort de la scène. Au dernier lancer.
\end{rpgnote}

\RPGEngravedPanel{Effort hors du combat}{assets/panel-17-engraved.png}{L'Érudit doit traduire une ancienne inscription avant l'effondrement du temple. Savoir Expert, SD 5. Il déclare un Effort avant le jet : 3d6, garder le plus élevé. Il obtient 2, 4 et 6. Réussite ! Après le jet, il paie 1F et tombe à 2 PV : Amoché. Quand le sol cède, il devra courir sur des jambes tremblantes.}

\RPGChapter{Guide du Narrateur}
\textbf{Le Seuil de difficulté est votre outil principal.} Le SD de base est 4. Augmentez-le pour les actions difficiles, réduisez-le pour les actions simples. Rien d'autre n'est nécessaire pour calibrer un défi.

\textbf{L'Épuisement engendre l'aventure.} Après l'échec d'une action Critique, le groupe ne peut pas simplement recommencer par la force : il doit trouver une autre voie. Cette contrainte encourage la créativité.

\textbf{Chaque ennemi veut quelque chose.} Un brigand veut de l'argent, pas nécessairement la mort. Un démon veut un pacte, pas une bagarre. Les motivations comptent davantage que les statistiques.

\textbf{Demandez toujours « Comment fais-tu ? »} La réponse détermine la Caractéristique utilisée et le SD.

\textbf{Interprétez le degré de réussite.} Un 6 face à SD 3 n'équivaut pas à un 4 face à SD 4. Les deux réussissent, mais le premier brillamment et le second de justesse. Jouez cette nuance.

\textbf{Les Blessures de Fatigue sont un outil.} Toute action extrême peut en coûter : soulever un rocher, courir des heures, résister à un poison. Elles ne concernent pas uniquement la magie.

\textbf{La magie est une décision collective.} Le Contrecoup d'un Grimoire peut affecter tout le monde. Avant un lancement risqué, laissez le groupe en discuter.

\textbf{Les interactions sociales utilisent les mêmes règles.} Persuader, mentir, intimider : lancez Savoir + la Compétence pertinente contre un SD. Aucun système distinct n'est nécessaire pour les conflits sociaux.

\RPGBackCoverFrench{AU DERNIER LANCER}
\end{document}
